\section{Introduction}
\label{sec:introduction}

Fact-checking over document corpora is a central task in applied NLP: given a
claim and a collection of source documents, a system must decide whether the
corpus supports the claim, refutes it, or provides no relevant evidence.
Current approaches fall into two failure modes. Retrieval-augmented generation
(RAG) systems retrieve supporting passages and feed them to a language model,
which can \emph{hallucinate} verdicts on claims the corpus does not speak
to~\cite{lewis2020rag}. Knowledge-graph pipelines that extract structured
triples achieve better grounding but require rigid, hand-curated schemas and
cannot express that a question simply lacks an answer in the available
evidence~\cite{hogan2021kg}.

Neither class of system handles \emph{multi-hop} reasoning with
\emph{calibrated uncertainty}: the ability to chain intermediate facts across
documents while simultaneously signalling when the chain is absent or
incomplete. This gap is consequential. A compliance analyst who asks whether
company A indirectly supplies component B via an intermediate supplier needs
both the chain and an honest ``the corpus does not say'' when the chain is
missing.

\paragraph{This paper.}
We present \textbf{Infon}, a deployable knowledge graph reasoner that
fills this gap. Infon is not a neural language model; it is a symbolic
reasoner augmented by a learned prior. Its design is guided by five
engineering properties:

\begin{enumerate}
  \item \textbf{Cassette-native storage.}  Documents are stored as immutable,
    content-addressed \texttt{.inf} files that can live on S3 or a local
    filesystem via \texttt{fsspec}. Delta ingests append; nothing is rewritten.
    A manifest pruner skips cassettes that cannot contain the query anchor,
    cutting Parquet opens by 7--16$\times$ at scale and 278$\times$ on cache
    misses.

  \item \textbf{Schema-free operation.}  Plain-English questions are encoded
    with SPLADE-tiny~\cite{formal2021splade} (17~MB, Apache~2.0, no GPU)
    to an anchor vocabulary. Connective predicates are auto-inferred from the
    data: no ontology annotation is required. A Strands-powered \texttt{Analyst}
    agent routes natural-language questions to the appropriate retrieval
    primitive, citing every source it uses.

  \item \textbf{Dempster--Shafer mass with first-class} $\boldsymbol{m(\Theta)}$.
    Every verdict is a four-part mass $(m_S, m_R, m_U, m_\Theta)$ over the
    frame $\Omega = \{S, R, U\}$ (Supports, Refutes, Uncertain), where
    $m_\Theta = m(\Omega)$ is the vacuous element that rises to $1.0$ on
    unsupported claims rather than a confident wrong answer~\cite{shafer1976mathematical}.

  \item \textbf{MCTS and sheaf GNN multi-hop reasoning.}  A Monte Carlo Tree
    Search traverses the infon hypergraph, with a 140k-parameter frozen sheaf
    graph neural network~\cite{bodnar2022sheaf} serving as the terminal scorer.
    The GNN provides the MCTS prior that guides search toward high-confidence
    chains.

  \item \textbf{Schema evolution without reingestion.}  A categorical
    \texttt{SchemaFunctor} rewrites existing cassettes under a new ontology
    via a Kan pushforward. Old snapshots remain queryable; the migration is
    62$\times$ faster than re-extraction.
\end{enumerate}

\paragraph{Roadmap.}
Section~\ref{sec:background} establishes the formal background: Dempster--Shafer
theory, situation semantics and infons, SPLADE sparse retrieval, and sheaf
graph neural networks.
Section~\ref{sec:architecture} describes Infon's four-layer architecture
(cassette substrate, query DSL, reasoner, and conversational analyst).
Section~\ref{sec:implementation} details the implementation, including the
synthgen corpus, GNN training, and manifest pruner benchmarks.
Section~\ref{sec:experiments} reports experimental results on synthetic and
real-world evaluations (HoVer~\cite{jiang2020hover},
AVeriTeC~\cite{schlichtkrull2023averitec},
SciFact~\cite{wadden2020scifact}).
Section~\ref{sec:discussion} discusses limitations and future work.
Section~\ref{sec:related_work} surveys related work.
Section~\ref{sec:conclusion} concludes.
Section~\ref{sec:reproducibility} provides a system card and reproducibility
statement.
